现在设 $\mE_1$ 与 $\mE_2$ 是 $D_R$ 上两个 $G$-torsor，$\beta: \mE_1\dashrightarrow \mE_2$ 为修正。
则对每个几何点 $x\in \Spec R$，经基变换得到三元组 $(\mE_1|_{D_x},\mE_2|_{D_x},\beta_x: \mE_1|_{D_x}\dashrightarrow \mE_2|_{D_x})$。令 $\inv_x(\beta):=\inv(\beta_x)$。

我们如下定义（球面）Schubert 胞腔 $\mathring{\Gr}_{\mu}\subset \Gr\otimes \bar k$：
\[\mathring{\Gr}_\mu:=\left\{(\mE,\beta)\in \Gr\otimes \bar k\mid \inv(\beta)=\mu\right\}.\]
可以证明它在 $\Gr\otimes \bar k$ 中局部闭，并且是过 $\mu(\varpi)$ 的 $L^+G\otimes \bar k$-轨道。其闭包记作 $\Gr_\mu$，称为（球面）Schubert 簇\footnote{注意文献中有时把 Schubert 簇记作 $\Gr_{\leq\mu}$ 或 $\overline{\Gr}_\mu$，而把 Schubert 胞腔记作 $\Gr_\mu$。希望我们的记号不至于太令人困惑。}。它是一些 Schubert 胞腔之并。于是可在 $\xcoch(T)/W$ 上定义偏序（通常称为 Bruhat 序）如下：
\[\la\leq \mu \quad \mbox{ 当且仅当 } \quad \Gr_{\la}\subset \Gr_\mu.\]
\begin{rmk}
通常，Bruhat 序用 $G$ 对应的根数据以组合方式定义。对我们后文的讨论，上述定义已经够用。
\end{rmk}
\begin{ex}
设 $G=\GL_n$。如例 \ref{invbetaGLn} 那样把 $\xcoch(T)/W$ 与 $\xcoch(T)^+$ 等同。
则 Bruhat 序可显式描述如下：$\mu_1=(m_1,\ldots,m_n)\leq \mu_2=(l_1,\ldots,l_n)$ 当
$$m_1+\cdots+m_j\leq l_1+\cdots+l_j, \ \ \ j=1,\ldots, n, \quad \mbox{ 且 } \ \ m_1+\cdots+m_n=l_1+\cdots+l_n.$$ 
\end{ex}

Schubert 簇 $\Gr_\mu$ 一般是定义在 $\bar k$ 上的射影簇\footnote{事实上，它定义在 $\mu$ 的定义域上。}。它完全光滑，当且仅当 $\Gr_\mu=\mathring{\Gr}_\mu$；此时 $\mu$ 是 $\xcoch(T)/W$ 中关于 Bruhat 序的极小元。这样的 $\mu$ 称为\textbf{极小}（minuscule）。

\begin{ex}
设 $G=\GL_n$。令 $\omega_i=1^i0^{n-i}$ 如例 \ref{minu Sch}，并令 $\omega_i^*=\omega_{n-i}-\omega_{n}$。
注意 $\inv(\beta)=\omega_i$ 当且仅当 $\beta$ 延拓为真正的映射 $\La_1\to \La_2$，使得 $\varpi \La_2\subset \La_1$，且 $\La_2/\La_1$ 是 $i$ 维 $k$-向量空间。类似地，$\inv(\beta)=\omega_i^*$ 当且仅当 $\beta^{-1}$ 诱导包含 $\varpi \La_1\subset \La_2\subset \La_1$，且 $\La_1/\La_2$ 的维数为 $i$。
不难看出 $\omega_i,\omega_i^*$ 是极小余权，且 $\Gr_{\omega_i}\cong \Gr_{\omega_{n-i}^*}\cong \Gr(n-i,n)$。取 $k$-点即回到例 \ref{minu Sch}。
\end{ex}

接下来我们简要讨论几何 Satake 的通常表述，更详细的讨论参看 \cite{Zh2}（至少在等特征情形）。
动机版本将在下一小节讨论。
为简单起见，本小节余下部分假定 $k=\bar{k}$，并把 $L^+G$ 记作 $K$、把 $L^nG$ 记作 $K_n$（于是与前面的记号相反，本小节余下部分中 $K$ 是投射代数群而非投射有限群）。令 $K^{(n)}=\ker(K\to K_n)$ 表示第 $n$ 级同余子群。我们固定系数环 $E$，使 $\on{char} k$ 在其中可逆。层均指 $E$-层。

我们先简要回忆仿射 Grassmann 簇上反常层范畴 $\on{P}_{K}(\Gr)$ 的定义，它通常称为\textbf{Satake 范畴}，记作 $\Sat_G$。
首先，可以选取表示 $\Gr=\underrightarrow\lim X_i$，使每个 $X_i$ 是 $K$-不变闭子概形（例如取 $X_i$ 为若干 Schubert 簇的有限并）。此外，$K$ 在 $X_i$ 上的作用经某个 $K_m$ 分解。
于是可定义
\[\Sat_G=\on{P}_{K}(\Gr)=\underrightarrow\lim\on{P}_{K}(X_i),\]
其中 $\on{P}_{K}(X_i)$ 定义为 $X_i$ 上 $K_m$-等变反常层的范畴（它与 $m$ 的选取无关，相差典范等价），连接函子 $\on{P}_{K}(X_i)\to \on{P}_{K}(X_{i'})$ 为沿闭浸入 $X_i\to X_{i'}$ 的前推。此范畴定义的详细讨论见 \cite[\S A.1.17]{Zh2}。

可以证明它是半单 Abel 幺半范畴，幺半结构由 Lusztig 的层卷积积给出。我们回忆其定义。令 $\Gr\tilde\times \Gr$ 为仿射 Grassmann 簇的扭积（也称卷积 Grassmann 簇）。就模空间问题而言，它定义为预层
\begin{equation*}\label{I:conv}
(\Gr\tilde\times\Gr)(R)=\left\{(\mE_1,\mE_2,\beta_1,\beta_2)\ \left|\ \begin{split}&\mE_1,\mE_2 \mbox{ 是 }  D_R\mbox{ 上的 } G\mbox{-torsor}, \\
&  \beta_1:\mE_1\dashrightarrow \mE^0, \beta_2:\mE_2\dashrightarrow  \mE_1\end{split}\right.\right\}.
\end{equation*}
存在自然映射
\begin{equation}\label{conv m0}
m:\Gr\tilde\times\Gr\to \Gr,\quad (\mE_1,\mE_2,\beta_1,\beta_2)\mapsto (\mE_2, \beta_1\beta_2).
\end{equation}
与自然投影
$$\pr_1:\Gr\tilde\times\Gr\to \Gr, \quad (\mE_1,\mE_2,\beta_1,\beta_2)\mapsto (\mE_1,\beta_1),$$
二者一起诱导同构
$(\pr_1,m):\Gr\tilde\times\Gr\cong \Gr\times \Gr$。
特别地，卷积 Grassmann 簇可表。映射 $m$ 通常称为卷积映射。

更一般地，存在 $r$ 重卷积 Grassmann 簇 $\Gr\tilde\times\cdots\tilde\times\Gr$，它分类 $G$-torsor 的修正链
\begin{equation}
\label{E:chain mod}
\mE_r\stackrel{\beta_r}{\dashrightarrow}\mE_{r-1}\stackrel{\beta_{r-1}}{\dashrightarrow}\cdots\stackrel{\beta_0}{\dashrightarrow}\mE_0=\mE^0.
\end{equation}
对每个 $j=1,\ldots,r-1$，存在到 $j$ 重卷积 Grassmann 簇的自然投影 $\pr_j$，即记住 $\mE_j\stackrel{\beta_j}{\dashrightarrow}\cdots\stackrel{\beta_0}{\dashrightarrow}\mE_0=\mE^0$；还存在 $r$ 重卷积映射
\begin{equation}\label{conv m}
m:\Gr\tilde\times\cdots\tilde\times\Gr\to \Gr,
\end{equation}
把 \eqref{E:chain mod} 送到 $(\mE_r, \beta_1\cdots\beta_r)$。如前，$K$ 通过 $\mE_0=\mE^0$ 的自同构自然作用在 $\Gr\tilde\times\cdots\tilde\times\Gr$ 上，且有 $K$-torsor $(\Gr\tilde\times\cdots\tilde\times\Gr)^{(\infty)}\to\Gr\tilde\times\cdots\tilde\times\Gr$，它分类 \eqref{E:chain mod} 连同同构 $\epsilon:\mE_r\simeq \mE^0$。
